\scanpage{200}
\mspage{\unsure{56}}
n'est pas forcément $\varprojlim$ de foncteurs représentables. \deleted{Donc} il est douteux que $c')$ suffise à entraîner que $u$ est $W_0$-connexe.

\uline{Corollaire} 4.2. Pour que $u$ soit $W_0$-connexe, il \uline{suffit} qu'il soit une $W_0$-équivalence universelle, i.e. \inserted{(cf \unsure{3.1|3.3})} que les fibres soient $0$-connexes, et que pour les flèches non identiques dans $I$, \unsure{définissant} $\unsure{\mathbb{\Delta}|\Delta}^1\mathbin{\unsure{-|\to}}I$, $I'\times_I\mathbb{\Delta}^1$ soit $0$-connexe.

4.3. Cette condition de $W_0$-équivalence universelle \deleted{\unsure{e}} est cependant \uline{strictement} plus forte que la $W_0$-connexité, comme on le voit sur l'exemple \deleted{d'une \unsure{équi}} où $u$ est une \uline{équivalence de catégories} non surjective\unsure{.} sur les objets (p. ex. l'inclusion d'un sommet dans un groupoïde $1$-connexe ayant plus de \unsure{un} objet). Plus généralement, tout foncteur de localisation est $W_0$-connexe, alors que ses fibres ne sont pas forcément connexes, \deleted{\unsure{et}} a fortiori il n'est pas une $W_0$-\unsure{équiv.} universelle.

Cependant, si \deleted{dans un $I$} tout \unsure{endomorphisme|automorphisme} $p$ d'un $i$ de $I$, tel que $p^2=p$\unsure{,|.} (projecteur) est l'identité, et si
% archive sheet number: 200
